\section{Related Work}

We review uncertainty-based hallucination detection methods and their evaluation paradigms, highlighting the absence of systematic cross-benchmark transfer analysis that motivates our work.

\subsection{Uncertainty Quantification for LLMs}

Semantic entropy \cite{kuhn2023semantic} computes uncertainty over clusters of semantically equivalent generations, achieving approximately 0.85 AUROC on TriviaQA. By using bidirectional entailment to group generations by meaning rather than surface form, semantic entropy captures linguistic invariances that token-level entropy misses. However, this work---like most in the field---evaluates each benchmark independently without testing cross-benchmark transfer.

P(True) methods \cite{kadavath2022language} prompt models to predict the probability that their own outputs are correct. Calibration improves with model scale, but the approach requires explicit self-evaluation prompting and has not been evaluated across diverse benchmark families. Our work focuses on semantic entropy, which operates during generation rather than requiring post-hoc evaluation.

Token-level uncertainty measures including entropy and predictive variance have been explored for uncertainty estimation \cite{xiao2021hallucination}, but semantic entropy consistently outperforms these approaches by accounting for meaning-level rather than surface-level variation.

\subsection{Consistency-Based Detection}

SelfCheckGPT \cite{manakul2023selfcheckgpt} detects hallucinations by measuring consistency across multiple samples, achieving strong performance on WikiBio without requiring external knowledge. While methodologically distinct from entropy-based approaches, consistency methods face similar transfer questions: calibration on one domain does not guarantee performance on another. Our clustering framework could extend to consistency-based methods in future work.

\subsection{Benchmark Evaluation Paradigms}

Existing hallucination benchmarks---TriviaQA \cite{joshi2017triviaqa}, HaluEval \cite{li2023halueval}, FEVER \cite{thorne2018fever}, and others---have been treated as interchangeable representatives of ``factual QA.'' However, no prior work has examined whether benchmarks cluster into distinct families based on the error processes they probe.

Domain adaptation literature \cite{bendavid2010theory} provides theoretical grounding for our findings: distribution shift degrades transfer, and distribution similarity predicts transferability. We operationalize this theory for hallucination detection, using JS-divergence to measure uncertainty distribution similarity and hierarchical clustering to discover benchmark families.
